% Generated from validated Article Draft Result. Do not hand-edit; revise the structured draft instead.
\section{言語モデルの外へ――CLIPとDALL-Eが開いた視覚との接続}
\label{pkg:multimodal-substrate}
\noindent\textbf{CLIPとDALL-Eは同じ『画像AI』ではない。CLIPはvision-language representationを、DALL-Eはtext-to-image generationを代表し、2021年初頭には表現を接続する層と画像を生成する層が別々の技術問題として現れていた。この分離を保つことが、その後のmultimodal generative AIを理解する出発点になる。}\autocite{src-005ad180fc46,src-009747d6fa68,src-bc37d5969edc,src-f14c51ae873c}

CLIPが2021年の記録で担う役割は、画像を生成することではなく、visionとlanguageをcontrastiveなrepresentationとして接続する系統を可視化した点にある。言語と視覚を同じ問題設定の中で扱えるrepresentation substrateが現れたことで、textだけを閉じた入力空間として扱う見方から、異なるmodalitiesを対応付ける設計へ視野が広がった。本稿ではCLIPの個別benchmark上の優位性を独立再現したとは扱わず、このrepresentation-levelの境界を採用する。\autocite{src-005ad180fc46,src-bc37d5969edc}


DALL-Eは別の枝を担う。textからimageを生成するgenerative modelingは、vision-language representationを作ることとは目的も出力も異なる。2021年に両者が近接して現れたことは、language model周辺の技術が視覚へ拡張したという一語では圧縮できない。『意味空間を接続する』ことと『新しい視覚artifactを生成する』ことを分けることで、multimodal substrateとgenerative mediaの二層構造が見える。\autocite{src-009747d6fa68,src-f14c51ae873c}


この区別は後続の技術を読むときにも効く。representation substrateは、入力や意味の対応関係をどう持つかという基盤側の問題であり、generationはその条件から何を生成するかという生成側の問題である。両者は相互に関係し得るが、同一の機構でも同一の成果でもない。2021年を『text-to-imageが始まった年』だけとして読むより、視覚と言語を接続するための複数のlayerが同時に立ち上がった年と見る方が技術史として正確である。\autocite{src-005ad180fc46,src-009747d6fa68,src-bc37d5969edc,src-f14c51ae873c}


この二層構造は、同じ年に別系統として進むdiffusionを理解するための足場にもなる。DALL-Eの存在だけから後年の画像生成ecosystemを直線的に導くのではなく、2021年時点ではrepresentation、autoregressiveなtext-to-image generation、そして別章で扱うdiffusionという複数の系統が並行していた、と整理するのが適切である。後の統合を知っているからこそ、当時まだ分離していた技術境界を残す必要がある。\autocite{src-005ad180fc46,src-009747d6fa68,src-bc37d5969edc,src-f14c51ae873c}


\begin{claimboundary}
CLIPとDALL-Eの技術的・性能上の評価は原著者・project側の主張として保持する。本稿が一次資料から確認するのは、両者が2021年のmultimodal technical recordを構成し、representationとgenerationという異なる役割を担ったことまでである。\autocite{src-005ad180fc46,src-009747d6fa68,src-bc37d5969edc,src-f14c51ae873c}
\end{claimboundary}
